LiionTR formulates thermal runaway as a coupled system of ordinary
differential equations integrated in time.

The numerical solver owns the complete dynamic state and evaluates the
physical models only as functions of the current state. This avoids
hidden mutable state inside the thermal, reaction, gas, pressure, and
venting models.


\subsection{State vector}

For a problem containing $N_r$ reaction states and $N_s$ gas species,
the integrated state vector is

\begin{equation}
    \mathbf{y}
    =
    \begin{bmatrix}
        T &
        \alpha_1 &
        \cdots &
        \alpha_{N_r} &
        n_1 &
        \cdots &
        n_{N_s}
    \end{bmatrix}^{\mathrm{T}}.
    \label{eq:numerical_state_vector}
\end{equation}

The first component is the lumped cell temperature.

The following $N_r$ components contain reaction conversions,

\begin{equation}
    \boldsymbol{\alpha}
    =
    \begin{bmatrix}
        \alpha_1 &
        \cdots &
        \alpha_{N_r}
    \end{bmatrix}^{\mathrm{T}},
\end{equation}

while the final $N_s$ components contain the gas-species inventories,

\begin{equation}
    \mathbf{n}
    =
    \begin{bmatrix}
        n_1 &
        \cdots &
        n_{N_s}
    \end{bmatrix}^{\mathrm{T}}.
\end{equation}

When gas generation is disabled, the gas portion of the state vector is
absent.


\subsection{Initial conditions}

The initial numerical state is assembled from the configured thermal
problem.

The temperature is initialized as

\begin{equation}
    T(0)
    =
    T_0.
\end{equation}

Each reaction is initialized according to its prescribed initial
conversion,

\begin{equation}
    \alpha_i(0)
    =
    \alpha_{i,0}.
\end{equation}

When gas modeling is enabled, the initial gas inventory is determined
from the configured species inventory or from the pressure model.

The complete initial state is therefore

\begin{equation}
    \mathbf{y}(0)
    =
    \mathbf{y}_0.
\end{equation}


\subsection{Right-hand-side evaluation}

At every solver evaluation, the current state vector is unpacked into
temperature, reaction conversions, and gas inventories.

The reaction context is constructed from the current conversion vector,

\begin{equation}
    \boldsymbol{\alpha}(t).
\end{equation}

The reaction network then evaluates

\begin{equation}
    \dot{\boldsymbol{\alpha}}
    =
    \mathbf{f}_{\alpha}
    \left(
        T,
        \boldsymbol{\alpha}
    \right).
    \label{eq:rhs_reaction_rates}
\end{equation}

The corresponding total reaction heat is

\begin{equation}
    \dot{Q}_{\mathrm{gen}}
    =
    \dot{Q}_{\mathrm{gen}}
    \left(
        T,
        \boldsymbol{\alpha}
    \right).
\end{equation}

The thermal model evaluates

\begin{equation}
    \dot{T}
    =
    \frac{
        \dot{Q}_{\mathrm{gen}}
        -
        hA(T-T_{\infty})
    }{
        C_{\mathrm{th}}
    }.
    \label{eq:rhs_temperature}
\end{equation}


\subsection{Gas-generation contribution}

If a gas-generation model is present, the reaction rates are also used
to determine species-generation rates,

\begin{equation}
    \dot{\mathbf{n}}_{\mathrm{gen}}
    =
    \mathbf{f}_{\mathrm{gas}}
    \left(
        \dot{\boldsymbol{\alpha}}
    \right).
    \label{eq:rhs_gas_generation}
\end{equation}

Before vent opening,

\begin{equation}
    \dot{\mathbf{n}}
    =
    \dot{\mathbf{n}}_{\mathrm{gen}}.
    \label{eq:rhs_closed_gas}
\end{equation}


\subsection{Pressure reconstruction}

Internal pressure is not integrated as an independent state variable.

Instead, it is reconstructed algebraically from the instantaneous gas
inventory and temperature,

\begin{equation}
    P
    =
    P
    \left(
        T,
        \mathbf{n}
    \right).
    \label{eq:pressure_algebraic_state}
\end{equation}

For the current ideal-gas formulation,

\begin{equation}
    P
    =
    \frac{
        R T
        \displaystyle\sum_s n_s
    }{
        V_{\mathrm{free}}
    }.
\end{equation}

Avoiding pressure as an additional integrated variable ensures direct
consistency between pressure and gas inventory.


\subsection{Vent-opening event}

When venting is enabled, the closed-cell integration defines an event
function

\begin{equation}
    g_{\mathrm{vent}}
    \left(
        t,
        \mathbf{y}
    \right)
    =
    P
    \left(
        T,
        \mathbf{n}
    \right)
    -
    P_{\mathrm{vent}}.
    \label{eq:vent_event_function}
\end{equation}

Vent opening occurs when

\begin{equation}
    g_{\mathrm{vent}}
    =
    0.
\end{equation}

The event terminates the closed-cell integration.


\subsection{Two-phase integration}

The current solver represents irreversible vent opening using two
successive numerical integrations.

The first phase solves the closed-cell system,

\begin{equation}
    \dot{\mathbf{y}}
    =
    \mathbf{F}_{\mathrm{closed}}
    \left(
        t,
        \mathbf{y}
    \right),
\end{equation}

until either the final time or a terminating event is reached.

If the vent-opening event occurs, the state at the event time

\begin{equation}
    \mathbf{y}_{\mathrm{vent}}
\end{equation}

becomes the initial condition of the second phase.

The second integration solves

\begin{equation}
    \dot{\mathbf{y}}
    =
    \mathbf{F}_{\mathrm{open}}
    \left(
        t,
        \mathbf{y}
    \right),
\end{equation}

where the gas balance is modified to include vent discharge,

\begin{equation}
    \dot{\mathbf{n}}
    =
    \dot{\mathbf{n}}_{\mathrm{gen}}
    -
    \dot{\mathbf{n}}_{\mathrm{vent}}.
    \label{eq:rhs_open_gas}
\end{equation}


\subsection{Solver flow}

The numerical sequence is summarized in
\cref{fig:solver_flow}.

\begin{figure}[htbp]
    \centering

    \begin{tikzpicture}[
        node distance=1.25cm and 2.0cm,
        block/.style={
            rectangle,
            rounded corners,
            draw,
            align=center,
            minimum width=4.0cm,
            minimum height=0.9cm
        },
        decision/.style={
            diamond,
            draw,
            align=center,
            aspect=2.2,
            inner sep=1pt
        },
        line/.style={
            -{Latex[length=3mm]},
            thick
        }
    ]

    \node[block] (initial)
        {Construct initial state $\mathbf{y}_0$};

    \node[block, below=of initial] (closed)
        {Integrate closed-cell system};

    \node[decision, below=of closed] (ventevent)
        {Vent event?};

    \node[block, below left=1.5cm and 2.0cm of ventevent] (finish)
        {Assemble result};

    \node[block, below right=1.5cm and 2.0cm of ventevent] (open)
        {Restart from event state\\
         Integrate open-cell system};

    \node[block, below=of open] (merge)
        {Concatenate closed and open solutions};

    \draw[line] (initial) -- (closed);
    \draw[line] (closed) -- (ventevent);

    \draw[line]
        (ventevent)
        --
        node[above left] {No}
        (finish);

    \draw[line]
        (ventevent)
        --
        node[above right] {Yes}
        (open);

    \draw[line] (open) -- (merge);
    \draw[line] (merge.west) -| (finish.south);

    \end{tikzpicture}

    \caption{
        Numerical integration strategy used for irreversible vent
        opening in the current LiionTR solver.
    }
    \label{fig:solver_flow}
\end{figure}


\subsection{Maximum-temperature event}

A maximum permitted temperature may be specified as

\begin{equation}
    T_{\max}.
\end{equation}

The corresponding event function is

\begin{equation}
    g_T
    =
    T-T_{\max}.
    \label{eq:max_temperature_event}
\end{equation}

The integration terminates when the event crosses zero in the prescribed
direction.


\subsection{Maximum-pressure event}

Similarly, a maximum-pressure event is represented by

\begin{equation}
    g_P
    =
    P-P_{\max}.
    \label{eq:max_pressure_event}
\end{equation}

This event can terminate either the closed or open phase of the
simulation.


\subsection{Stiff numerical integration}

Thermal runaway kinetics frequently contain strongly temperature-sensitive
Arrhenius terms spanning several orders of magnitude.

The coupled equations can consequently become stiff, particularly during
the rapid acceleration preceding thermal runaway.

The default LiionTR solver therefore uses the backward differentiation
formula method available through SciPy's \texttt{solve\_ivp} interface
\parencite{virtanen2020scipy}.

The BDF method is an implicit variable-order integration method suitable
for stiff ordinary differential equations.


\subsection{Adaptive time stepping}

The solver automatically adjusts its internal time step according to
local error estimates.

The user may specify relative and absolute tolerances,

\begin{equation}
    \mathrm{rtol}
\end{equation}

and

\begin{equation}
    \mathrm{atol},
\end{equation}

as well as an optional maximum time step.

The internal time-step size can become substantially smaller during the
rapid thermal runaway phase than during the initial slow-heating period.


\subsection{State reconstruction}

After numerical integration, the resulting solver state is separated
into physically meaningful output quantities.

These include

\begin{itemize}
    \item time;
    \item cell temperature;
    \item individual reaction conversions;
    \item reaction progress rates;
    \item heat-generation rates;
    \item gas-species inventories;
    \item total gas amount;
    \item gas composition;
    \item pressure;
    \item vent status.
\end{itemize}

For vented simulations, the two numerical phases are concatenated into a
single continuous simulation result.


\subsection{Numerical architecture}

The numerical solver acts as the orchestration layer between the
individual physical models.

Conceptually,

\begin{equation}
    \mathbf{y}
    \longrightarrow
    \begin{cases}
        \text{reaction model},\\
        \text{thermal model},\\
        \text{gas model},\\
        \text{pressure model},\\
        \text{vent model}
    \end{cases}
    \longrightarrow
    \dot{\mathbf{y}}.
\end{equation}

Individual physical components therefore remain independent of the
specific numerical integration algorithm.


\subsection{Current limitations}

The current numerical formulation assumes:

\begin{itemize}
    \item a single lumped temperature state;
    \item scalar reaction conversions;
    \item species-resolved gas inventories;
    \item pressure reconstructed algebraically;
    \item instantaneous transition from closed to fully open vent;
    \item no resealing after vent opening;
    \item no differential-algebraic constraints beyond algebraic
          reconstruction of derived quantities;
    \item no spatial discretization.
\end{itemize}


\subsection{Future numerical extensions}

Future developments may include:

\begin{itemize}
    \item alternative stiff and non-stiff integrators;
    \item explicit Jacobian evaluation;
    \item sparse Jacobian support;
    \item differential-algebraic formulations;
    \item adaptive vent-opening models;
    \item integration of internal energy instead of temperature;
    \item spatial thermal discretization;
    \item cell-to-cell propagation systems;
    \item sensitivity equations and parameter estimation.
\end{itemize}